\setcounter{section}{20}

  \zwischenueberschrift{Übungsaufgaben}

\inputaufgabe
{}
{

Bestimme die
\definitionsverweis {äußere Ableitung}{}{}
der
$1$-\definitionsverweis {Differentialform}{}{}

\mavergleichskettedisp
{\vergleichskette
{ \omega
}
{ =} { \left(  x^2-y^3  \right) dx+x^3y^2dy
}
{ } {
}
{ } {
}
{ } {
}
}
{}{}{}
auf dem $\R^2$.

}
{} {}

\inputaufgabe
{}
{

Bestimme die
\definitionsverweis {äußere Ableitung}{}{}
der
$1$-\definitionsverweis {Differentialform}{}{}

\mavergleichskettedisp
{\vergleichskette
{ \omega
}
{ =} { xy^2dx+yzdy+x^3dz
}
{ } {
}
{ } {
}
{ } {
}
}
{}{}{}
auf dem $\R^3$.

}
{} {}

\inputaufgabegibtloesung
{}
{

Berechne die
\definitionsverweis {äußere Ableitung}{}{}
$d \omega$ der
$1$-\definitionsverweis {Differentialform}{}{}

\mavergleichskettedisp
{\vergleichskette
{ \omega
}
{ =} { { \frac{   x^2 }{  y } } dx- { \frac{   x  }{  y^2 } } dy
}
{ } {
}
{ } {
}
{ } {
}
}
{}{}{}
auf

\mavergleichskette
{\vergleichskette
{ U
}
{ = }{ \left\{ (x,y) \in \R^2  \mid   y \neq 0         \right\}
}
{  }{
}
{    }{
}
{   }{
}
}
{}{}{.}

}
{} {}

\inputaufgabegibtloesung
{}
{

Berechne die äußere Ableitung $d \omega$ der Differentialform

\mavergleichskettedisp
{\vergleichskette
{ \omega
}
{ =} {  e^{xz} dx \wedge dy - xyz dx \wedge dz + \left(  \sin ( \cos (xy) )  +y^{10}z^{100}  \right) dy \wedge dz
}
{ } {
}
{ } {
}
{ } {
}
}
{}{}{}
auf dem
\mathl{\R^3}{.}

}
{} {}

\inputaufgabegibtloesung
{}
{

Es sei
\maabbeledisp {f} {\R} {\R
} { t } {f(t)
} {,}
die durch

\mavergleichskettedisp
{\vergleichskette
{ f(t)
}
{ =} { { \frac{   \sin^{ 3  } (t^4)  }{  1+t^2 } }
}
{ } {
}
{ } {
}
{ } {
}
}
{}{}{}
gegeben ist.

a) Berechne die äußere Ableitung von $f$.

b) Berechne die äußere Ableitung von $fdt$.

}
{} {}

\inputaufgabe
{}
{

Bestimme die
\definitionsverweis {äußere Ableitung}{}{} der
$2$-\definitionsverweis {Differentialform}{}{}

\mathdisp {\omega = xdx \wedge dy+xy^2zdy \wedge dz+xe^ydx\wedge dz} {  }
 auf dem $\R^3$.

}
{} {}

\inputaufgabegibtloesung
{}
{

Zeige, dass die
$1$-\definitionsverweis {Differentialform}{}{}

\mavergleichskette
{\vergleichskette
{ \omega
}
{ = }{ (2x- \sin  y )dx-x\cos  ydy
}
{  }{
}
{    }{
}
{   }{
}
}
{}{}{}
auf dem $\R^2$
\definitionsverweis {geschlossen}{}{}
und auch
\definitionsverweis {exakt}{}{}
ist.

}
{} {}

\inputaufgabe
{}
{

Es sei

\mavergleichskette
{\vergleichskette
{ U
}
{ \subseteq }{ \R^n
}
{  }{
}
{    }{
}
{   }{
}
}
{}{}{}
eine
\definitionsverweis {offene Teilmenge}{}{}
und

\mavergleichskette
{\vergleichskette
{ \omega
}
{ = }{ \sum_{i = 1}^n f_i dx_i
}
{  }{
}
{    }{
}
{   }{
}
}
{}{}{}
eine
$1$-\definitionsverweis {Differentialform}{}{,}
die mit
\definitionsverweis {stetig differenzierbaren Funktionen}{}{}
\maabb {f_i} { U } { \R
} {}
beschrieben werde. Zeige, dass $\omega$ genau dann
\definitionsverweis {geschlossen}{}{}
ist, wenn

\mavergleichskettedisp
{\vergleichskette
{ \partial_i f_j
}
{ =} { \partial_j f_i
}
{ } {
}
{ } {
}
{ } {
}
}
{}{}{}
für alle $i,j$  gilt.

}
{} {}

\inputaufgabe
{}
{

Es sei

\mavergleichskette
{\vergleichskette
{ U
}
{ \subseteq }{ \R^n
}
{  }{
}
{    }{
}
{   }{
}
}
{}{}{}
eine
\definitionsverweis {offene Teilmenge}{}{}
und $\omega$ eine
\definitionsverweis {stetig differenzierbare}{}{}
$1$-\definitionsverweis {Form}{}{}
auf $U$ mit dem gemäß
[[Riemannsche Mannigfaltigkeit/Vektorfelder und 1-Formen/Fakt|Lemma 16.3]]
zugehörigen
\definitionsverweis {Vektorfeld}{}{}
$F$ auf $U$. Zeige, dass $\omega$ genau dann
\definitionsverweis {geschlossen}{}{}
ist, wenn $F$ die
\definitionsverweis {Integrabilitätsbedingung}{}{}
erfüllt.

}
{} {}

\inputaufgabe
{}
{

Es sei
\mathl{U \subseteq \R^n}{} eine
\definitionsverweis {offene Teilmenge}{}{}
und $\omega$ eine
\definitionsverweis {stetig differenzierbare}{}{}
$1$-\definitionsverweis {Form}{}{}
auf $U$ mit dem gemäß
[[Riemannsche Mannigfaltigkeit/Vektorfelder und 1-Formen/Fakt|Lemma 16.3]]
zugehörigen
\definitionsverweis {Vektorfeld}{}{}
$F$ auf $U$. Zeige, dass $\omega$ genau dann
\definitionsverweis {exakt}{}{}
ist, wenn $F$ ein
\definitionsverweis {Gradientenfeld}{}{}
ist.

}
{} {}

\inputaufgabe
{}
{

Es sei $\omega$ eine
\definitionsverweis {differenzierbare}{}{} \definitionsverweis {Differentialform}{}{}
auf einer Mannigfaltigkeit $M$ und
\maabb {\varphi} { L } { M
} {}
eine auf der Mannigfaltigkeit $L$ definierte differenzierbare Abbildung.

\aufzaehlungzwei {Es sei $\omega$
\definitionsverweis {exakt}{}{.}
Zeige, dass auch die zurückgezogene Differentialform $\varphi^*\omega$ exakt ist.
} {Es sei $\omega$
\definitionsverweis {geschlossen}{}{.}
Zeige, dass auch die
\definitionsverweis {zurückgezogene}{}{}
Differentialform $\varphi^*\omega$ geschlossen ist.
}

}
{} {}

Für die folgenden Aufgaben vergleiche
[[Vektorfeld/(x,y) nach (-y,x)/Nicht integrabel/Beispiel|Beispiel 58.7 (Analysis (Osnabrück 2021-2023))]]
und
[[Vektorfeld/Punktierte Ebene/(x,y) nach (-y,x) durch x^2+y^2/Integrabel nicht exakt/Beispiel|Beispiel 58.10 (Analysis (Osnabrück 2021-2023))]].

\inputaufgabe
{}
{

Zeige, dass die
$1$-\definitionsverweis {Differentialform}{}{}

\mathdisp {-ydx+xdy} {  }

auf dem $\R^2$ nicht
\definitionsverweis {geschlossen}{}{}
ist.

}
{} {}

\inputaufgabe
{}
{

Zeige, dass die
$1$-\definitionsverweis {Differentialform}{}{}

\mathdisp {- { \frac{   y  }{  x^2+y^2 } }  dx+ { \frac{   x  }{  x^2+y^2 } } dy} {  }

auf dem $\R^2 \setminus \{0\}$
\definitionsverweis {geschlossen}{}{,}
aber nicht
\definitionsverweis {exakt}{}{}
ist.

}
{} {}

\inputaufgabegibtloesung
{}
{

Zeige, dass jede stetige
$n$-\definitionsverweis {Differentialform}{}{}
$\omega$ auf einer offenen Menge

\mavergleichskette
{\vergleichskette
{U
}
{ \subseteq }{\R^n
}
{  }{
}
{    }{
}
{   }{
}
}
{}{}{}
\definitionsverweis {exakt}{}{}
ist.

}
{} {}

\inputaufgabe
{}
{

Es sei $M$ eine
\definitionsverweis {zusammenhängende}{}{}
\definitionsverweis {differenzierbare Mannigfaltigkeit}{}{}
und $\omega$ eine
\definitionsverweis {differenzierbare}{}{}
$1$-\definitionsverweis {Form}{}{}
auf $M$. Zeige, dass $\omega$ genau dann
\definitionsverweis {exakt}{}{}
ist, wenn für jeden
\definitionsverweis {stetig differenzierbaren}{}{}
Weg
\maabbdisp {\gamma} { [a,b] } { M
} {}
das
\definitionsverweis {Wegintegral}{}{}

\mathl{\int_\gamma \omega}{} nur von
\mathl{\gamma(a)}{} und
\mathl{\gamma(b)}{} abhängt.

}
{} {}

\inputaufgabe
{}
{

Es seien $M$ und $N$
\definitionsverweis {differenzierbare Mannigfaltigkeiten}{}{}
und
\maabbdisp {\varphi} { M } { N
} {}
eine
\definitionsverweis {differenzierbare Abbildung}{}{.}
Es sei $\omega$ eine
\definitionsverweis {differenzierbare}{}{}
$n$-\definitionsverweis {Form}{}{}
auf $N$, wobei $n$ die Dimension von $N$ sei. Zeige, dass $\varphi^* \omega$ eine
\definitionsverweis {geschlossene}{}{}
\definitionsverweis {Differentialform}{}{}
auf $M$ ist.

}
{} {}

\inputaufgabe
{}
{

Welche der folgenden Funktionen
\maabbdisp {} {\R_+} {\R
} {}
lassen sich
\definitionsverweis {differenzierbar}{}{}
in den
\definitionsverweis {Randpunkt}{}{}
$0$ fortsetzen?
\aufzaehlungsechs{
\mathl{x^3+ \sin^{ 3  }  x -e^{-x}}{,}
}{
\mathl{{ \frac{   1  }{  x  } }}{,}
}{
\mathl{\sin  { \frac{   1  }{  x  } }}{,}
}{
\mathl{x \sin  { \frac{   1  }{  x  } }}{,}
}{
\mathl{{ \frac{   e^{ { \frac{   1  }{  x  } } }  }{  x  } }}{,}
}{
\mathl{x^2 \sin  { \frac{   1  }{  x  } }}{.}
}

}
{} {}

\inputaufgabe
{}
{

Es sei
\mathl{H \subset \R^n}{} ein
\definitionsverweis {Halbraum}{}{.}
Es sei
\mathl{Q\in H}{} ein Punkt und $Q \in U \subseteq H$, wobei $U$ eine
\definitionsverweis {offene Teilmenge}{}{}
des $\R^n$ sei. Zeige, dass $Q$ kein Randpunkt von $H$ ist.

}
{} {}

\inputaufgabe
{}
{

Definiere die Begriffe \stichwort {Diffeomorphismus} {,} \stichwort {totales Differential} {} und \stichwort {höhere Ableitungen} {} für
\definitionsverweis {Halbräume}{}{}
\zusatzklammer {bzw. offene Teilmengen davon} {} {.}

}
{} {}

  \zwischenueberschrift{Aufgaben zum Abgeben}

\inputaufgabe
{3}
{

Bestimme die
\definitionsverweis {äußere Ableitung}{}{} der
$1$-\definitionsverweis {Differentialform}{}{}

\mathdisp {\omega = xy^2z^3dx+xyzdy+x^3yz^4dz} {  }
 auf dem $\R^3$.

}
{} {}

\inputaufgabe
{3}
{

Bestimme die
\definitionsverweis {äußere Ableitung}{}{} der
$2$-\definitionsverweis {Differentialform}{}{}

\mathdisp {\omega = xy^2dx \wedge dy+(x^3-y^2z^4)dy \wedge dz+ \sin (xy) dx \wedge dz} {  }

auf dem $\R^3$.

}
{} {}

\inputaufgabe
{5}
{

Es sei

\mavergleichskette
{\vergleichskette
{ U
}
{ \subseteq }{ \R^n
}
{  }{
}
{    }{
}
{   }{
}
}
{}{}{}
\definitionsverweis {offen}{}{}
und es seien
\mathl{\omega_1 , \ldots , \omega_r}{} Differentialformen auf $U$, wobei $\omega_i$ eine
$k_i$-\definitionsverweis { Differentialform}{}{}
sei. Finde und beweise eine Formel für

\mathdisp {d( \omega_1 \wedge \ldots \wedge \omega_r)} { . }

}
{} {}

\inputaufgabe
{5}
{

Zeige, dass die
\definitionsverweis {Differentialform}{}{}

\mavergleichskettedisp
{\vergleichskette
{ \omega
}
{ =} { \left(  2xy+3x^2-y e^{xy}  \right) dx + \left(  x^2-xe^{xy} +8y  \right)  dy
}
{ } {
}
{ } {
}
{ } {
}
}
{}{}{}
auf dem $\R^2$
\definitionsverweis {geschlossen}{}{}
und auch
\definitionsverweis {exakt}{}{}
ist.

}
{} {}

\inputaufgabe
{6}
{

Zeige, dass die
\definitionsverweis {Halbebene}{}{}
$\R_{\geq 0} \times \R$ und der
\definitionsverweis {Quadrant}{}{}

\mathl{\R_{\geq 0} \times \R_{\geq 0}}{}
\definitionsverweis {homöomorph}{}{}
sind.

}
{} {}

 [[Kategorie:Latexseite]]
